\section{Operational Experience}
\label{sec:operations}

A platform operated by a single administrator is exercised primarily by incidents: failures that must be detected, diagnosed, and remediated without a dedicated operations team. This section documents one representative production incident as a case study of the operability objectives formulated in Chapter~\ref{cha:system-architecture-design}. The incident is instructive because it originates below the Kubernetes abstraction, in host-level package management on a worker node, yet manifests as a user-facing workload failure. It therefore delimits what container orchestration does and does not isolate.

\subsection{Case Study: GPU Driver Version Mismatch}
\label{sec:driver-mismatch}

On 22~August~2026, a user-requested GPU notebook crash-looped at container creation. The pod event stream contained the decisive error, \texttt{nvidia-container-cli: initialization error: nvml error: driver/library version mismatch}, in the NVIDIA runtime's pre-start hook. Detection was by user impact rather than by the platform: the cluster had no node-level GPU health monitoring at the time, so no alert preceded the failed spawn. The event stream alone sufficed to localize the fault to the GPU stack of one worker node and to exclude the scheduler, the storage layer, and the notebook image.

Running \texttt{nvidia-smi} on the node produced the identical error, ruling out the container layer. The node, up for 38~days, was running kernel 6.8.0-134 with the NVIDIA kernel module 580.159.03 loaded at boot, whereas Ubuntu's unattended-upgrades had silently updated the userspace driver libraries to 580.173.02 on the previous morning and installed a new kernel (6.8.0-138) with matching precompiled modules. NVML refuses to operate across this version skew, so every subsequent GPU container creation on the node failed. A module reload could not have restored service: the on-disk module for the running kernel was still 580.159.03, and a matching module existed only for the not-yet-booted kernel. A reboot into the new kernel was therefore the minimal correct remediation.

After the reboot, the kernel module and the userspace libraries agreed at 580.173.02, and both GPUs (an NVIDIA TITAN~V and a TITAN~Xp) reported healthy. One secondary fault remained: the notebook pod scheduled during the reboot window was stuck in a terminal \texttt{UnexpectedAdmissionError} state, rejected with ``Allocate failed due to no healthy devices present; cannot allocate unhealthy devices \texttt{volcano.sh/vgpu-memory}''. This was a stale admission decision made before the vGPU device plugin had re-registered its devices with the kubelet; deleting the pod allowed JupyterHub to spawn a healthy replacement. User-visible recovery from the first failed spawn took approximately seven minutes.

Two preventive measures follow. First, GPU driver upgrades must be coupled to a maintenance action: either reboot promptly after a driver upgrade, or hold \texttt{nvidia-*} and \texttt{linux-modules-nvidia-*} back from unattended-upgrades and upgrade only during planned windows. The version-skew window is otherwise unbounded; in this case it lasted 26~hours. Second, node-level GPU health checks, such as a periodic \texttt{nvidia-smi} probe or a DCGM exporter per GPU node, would convert such outages from user-visible to administrator-visible and shrink detection time independently of workload activity.

\subsection{Agentic AI in the Operations Loop}
\label{sec:agentic-ops}

This incident was diagnosed and remediated with the assistance of an agentic AI assistant granted tool access to the cluster: a read-only Kubernetes API and SSH access to nodes. The assistant autonomously executed the complete diagnostic chain, from pod event triage and node inspection through kernel-module and library version comparison, package-manager history, and a cross-kernel \texttt{modinfo} check that established the insufficiency of a module reload. It then narrowed the response space to a small set of remediation options with explicit trade-offs: drain-and-reboot, plain reboot, or deferral to a maintenance window. The administrator's entire contribution was a single executive decision among these options, which the assistant executed and verified, including the resolution of the secondary admission failure and a post-recovery GPU smoke test.

This division of labor, with machine-executed diagnostics and human-owned judgment and accountability, directly supports the single-administrator operability objective of Chapter~\ref{cha:system-architecture-design}: it reduces the expertise and wall-clock time required per incident while keeping the administrator in the loop at the point where responsibility lies. The observed bottleneck was decision latency rather than diagnostic skill, which suggests that the marginal investment is not further automation of actions but better surfacing of options; the observability stack of Section~\ref{sec:observability} operationalizes exactly this conclusion.

